\begin{tikzpicture}[
  x=1mm,y=1mm,
  every node/.style={font=\small,text=BookInk},
  state/.style={circle,draw=BookTeal,line width=0.8pt,
    minimum size=8mm,inner sep=1pt},
  observed/.style={state,fill=FigureInput!18},
  indicator/.style={draw=BookAmber,fill=BookPaper,line width=0.8pt,
    minimum width=11mm,minimum height=8mm,inner sep=1pt}
]
  \begin{scope}
    \node at (12,12) {MCAR};
    \node[state] (h1) at (0,0) {$H$};
    \node[observed] (a1) at (24,0) {$A$};
    \node[indicator] (r1) at (12,-20) {$R_H$};
    \draw[book arrow] (h1) -- (a1);
    \node[align=center,text width=31mm] at (12,-34)
      {记录过程\\独立于数据值};
  \end{scope}
  \begin{scope}[xshift=35mm]
    \node at (12,12) {MAR};
    \node[state] (h2) at (0,0) {$H$};
    \node[observed] (a2) at (24,0) {$A$};
    \node[indicator] (r2) at (12,-20) {$R_H$};
    \draw[book arrow] (h2) -- (a2);
    \draw[book arrow] (a2) -- (r2);
    \node[align=center,text width=31mm] at (12,-34)
      {记录过程可依赖\\已知的报警};
  \end{scope}
  \begin{scope}[xshift=70mm]
    \node at (12,12) {MNAR};
    \node[state] (h3) at (0,0) {$H$};
    \node[observed] (a3) at (24,0) {$A$};
    \node[indicator] (r3) at (12,-20) {$R_H$};
    \draw[book arrow] (h3) -- (a3);
    \draw[book arrow] (a3) -- (r3);
    \draw[book arrow,draw=BookAmber,dashed] (h3) -- (r3);
    \node[align=center,text width=31mm] at (12,-34)
      {给定报警后仍\\依赖过热状态};
  \end{scope}
\end{tikzpicture}
